\subsection{Analytical Function}

\subsubsection{RealundefinedDerivativeundefinedandundefinedComplexundefinedDerivative}

Aundefinedcomplexundefinedfunctionundefined$f:\mathbb{C}\to \mathbb{C}$undefinedresembles,undefinedmoreundefinedorundefinedless,undefinedfunctionsundefinedlikeundefined$\mathbb{R}^2\to \mathbb{R}^2$.undefinedForundefinedinstance,undefinedtheundefinedcomplexundefinedexponentialundefined$f(z)=\exp(z)$undefinedwouldundefinedbeundefinedmoreundefinedorundefinedlessundefinedsimilarundefinedto
\begin{equation}
f\begin{pmatrix}
        x \\
        y
    \end{pmatrix} = 
    \begin{pmatrix}
        \exp(x)\cos(y) \\
        \exp(x)\sin(y)
    \end{pmatrix}
\end{equation}

Sinceundefinedtheundefinedcomplexundefinedversionundefinedcanundefinedbeundefinedrewrittenundefinedasundefined$\exp(z)=\exp(x+iy)=\exp(x)(\cos(y) + i\sin(y))$.

However,undefinedthereundefineddoesundefinedexistundefinedcrucialundefineddifferencies.undefinedRecallundefinedforundefinedrealundefinedfunctionsundefined$(\mathbb{R}^2, \mathbb{R}^2)$,undefinedtheyundefinedareundefineddifferentiableundefinedifundefinedthereundefinedexistsundefinedaundefinedlinearundefinedtransformationundefined$A\in \mathcal L (\mathbb{R}^2, \mathbb{R}^2)$undefinedsuchundefinedthat:
\begin{equation}
\lim_{x\to x_0} \frac{\|f(x_0+h)-f(x_0) - Ah\|}{\|h\|}=0
\end{equation}

Ifundefinedsuchundefined$A$undefinedexistsundefinedthenundefineditundefinedmustundefinedequalundefinedtoundefinedtoundefinedtheundefined\textbf{Jacobian}undefined$Df(x_0)$undefineddefinedundefinedas:
\begin{equation}
\left(\frac{\partial f_i}{\partial x_j}(x_0)\right)_{ij} \in \mathbb{R}^{m\times n}
\end{equation}

Furthermore,undefinedtheundefinedsufficientundefinedconditionundefinedforundefinedthisundefinedlinearundefinedtransforamtionundefined$A$undefinedtoundefinedexistsundefinedisundefinedgivenundefinedbyundefinedtheundefinedfollowingundefinedtheorem:

\begin{mystthmtheorem}{}{theorem-analyticalfunction-1}Ifundefined$f:\mathbb{R}^n\to \mathbb{R}^m$,undefinedthenundefined$Df(a)$undefinedeixstsundefinediffundefinedallundefined$D_jf^{i}(x_0)$undefinedexistundefinedinundefinedanundefinedopenundefinedsetundefinedcontainingundefined$x_0$undefinedandundefinedifundefinedeachundefinedfunctionundefined$D_jf^i$undefinedisundefinedcontinousundefinedatundefined$x_0$.
Rmk.undefinedweundefinedcallundefinedsuchundefined$f$undefined\textbf{continouslyundefineddifferentiable}

\end{mystthmtheorem}
Complexundefinedderivative,undefinedonundefinedtheundefinedotherundefinedhand,undefinedisundefineddefinedundefinedas:

\begin{mystthmdefinition}{}{definition-analyticalfunction-2}Letundefined$f:\mathbb{C} \to \mathbb{C}$undefinedbeundefinedaundefinedcomplexundefinedfunction.undefinedTheundefined\textbf{derivative}undefinedofundefined$f$undefinedatundefined$z_0$undefinedisundefineddefinedundefinedas
\begin{equation}
\frac{df}{dz}({z_0})= \lim_{\Delta z\to 0}\frac{f(z_0 + \Delta z) -f(z_0)}{\Delta z}
\end{equation}
Or,
\begin{equation}
\frac{df}{dz}({z_0})= \lim_{\Delta z\to 0}\frac{u(z+\Delta z) -u(z)}{\Delta z} + \lim_{\Delta z\to 0}\frac{v(z+\Delta z) -v(z)}{\Delta z} i
\end{equation}
whereundefined$u : \mathbb{C} \to \mathbb{R}$undefinedisundefinedtheundefinedrealundefinedpartundefinedofundefinedtheundefinedcomplexundefinedfunctionundefinedandundefined$v$undefinedbeingundefinedtheundefinedcomplexundefinedpart.

\end{mystthmdefinition}
Thisundefinedisundefined\textit{veryundefineddifferent}undefinedfromundefined$f:\mathbb{R}^2 \to \mathbb{R}^2$!undefinedForundefinedtheundefinedfollowingundefinedtwoundefinedreasons:

\begin{mystadmonition}{mystRed}{Differenceundefined1:undefinedDivision}\label{division}Ifundefinedweundefinedinsistundefinedwriteundefinedtheundefinedsameundefineddefinitionundefinedforundefinedtheundefinedrealundefinedanalog(i.e.undefineddenoteundefinedaundefinedcomplexundefinednumberundefined$a+bi$undefinedasundefined$a \choose b$,undefined$f_1\leftrightarrow u, f_2\leftrightarrow v$),undefinedthenundefinedtheundefinedaboveundefineddefinitionundefinedwouldundefinedlookundefinedlike:
\begin{equation}
\lim_{(x,y)^T \to (x_0,y_0)^T}
\begin{pmatrix}
\frac{\partial f_1}{\partial (x, y)} \\
\frac{\partial f_2}{\partial (x, y)}
\end{pmatrix}
\end{equation}
Whichundefineddoesn'tundefinedmakeundefinedsenseundefinedatundefinedall(Althoughundefinedoneundefinedcouldundefinedargueundefined$\frac{\partial f_i}{\partial (x,y)}$undefinedisundefinedjustundefined$(\nabla f_i)^T$undefinedbutundefinedthat'sundefinedmerelyundefinedsymbolic.undefinedWhatundefinedwouldundefineddivisionundefinedbyundefinedaundefinedvectorundefinedmean?).

\end{mystadmonition}
\begin{mystadmonition}{mystRed}{Differenceundefined2undefined:undefinedMixingundefinedOfundefinedComponents}\label{mixing-of-components}Toundefinedmakeundefineditundefinedworse,undefinedevenundefinedifundefinedweundefinedpretendundefinedsuchundefinednotationundefinedisundefinedsensible,undefinedtheyundefinedareundefinedstillundefineddifferent.undefinedTheundefinedlimitundefinedofundefinedtheundefinedrealundefinedpartundefined$\lim_{\Delta z\to 0}\frac{u(z+\Delta z) -u(z)}{\Delta z}$(whichundefinedweundefinedhopedundefinedtoundefinedcorrespondundefinedtoundefined$\frac{\partial f_1}{\partial (x,y)}$)undefinedhasundefinedaundefinedcomplexundefinedpartundefinedandundefinedsimilarilyundefinedtheundefinedlimitundefinedofundefinedtheundefinedcomplexundefinedpartundefinedalsoundefinedhasundefinedaundefinedrealundefinedpart!undefinedThisundefinedmeansundefinedtheundefinedcomponentsundefinedareundefinednaturallyundefinedmixing.

\end{mystadmonition}
Soundefinedtheundefinedrealundefineddifferenceundefinedisundefinedthatundefinedcomplexundefinednumberundefinedshouldundefinedbeundefinedviewedundefinedasundefinedaundefinedwholeundefinedwhereundefineddivisionundefinedisundefinedwell-defined.undefinedForundefinedme,undefinedifundefinedIundefinedeverundefinedfoundundefineditundefinedhardundefinedtoundefinedconvinceundefinedmyselfundefined$\mathbb{R}^2 \to \mathbb{R}^2$undefinedisundefineddifferentundefinedfromundefined$\mathbb{C} \to \mathbb{C}$,undefinedIundefinedwouldundefinedtryundefinedtoundefinedanswerundefinedtheundefinedquestionundefined``whyundefinedcan'tundefinedweundefineduseundefined$(a,b)$undefinedtoundefineddenoteundefineda+bi''?.undefinedThisundefinedtendundefinedtoundefinedclearundefinedthingsundefinedup.

Theundefinednextundefinedexampleundefinedillustrateundefinedthisundefinedkeyundefineddifference:

\begin{mystthmexample}{Aundefinednotundefineddifferentiableundefinedcomplexundefinedfunctionundefinedwithundefineditsundefinedrealundefinedanalogundefineddifferentiable}{example-analyticalfunction-3}Considerundefined$f(x+iy)=x^2 + iy$.undefinedTheundefinedrealundefinedanalog(InteractiveundefinedPlotundefinedOfundefinedTheundefinedRealundefinedAnalog)undefinedis:
\begin{equation}
f\begin{pmatrix}
x \\
y
\end{pmatrix}= 
\begin{pmatrix}
x^2 \\
y
\end{pmatrix}
\end{equation}
Thisundefinedfunctionundefinedisundefinedclearlyundefineddifferentiableundefinedsinceundefinedtheundefinedjacobianundefinedis:
\begin{equation}
\begin{pmatrix}
2x & 0 \\
0 & 1
\end{pmatrix}
\end{equation}
whichundefinedclearlyundefinedcontinuous(component-wise)undefinedforundefinedallundefined$x,y$.

However,undefinedthisundefinedisundefinednotundefinedtrueundefinedforundefinedtheundefinedcomplexundefinedanalog.undefinedWeundefinednowundefinedshowundefinedtheundefinedcomplexundefinedfunctionundefinedisundefinednotundefineddifferentiableundefinedatundefined$0\in \mathbb{C}$.

Supposeundefineditundefinedis,undefinedthenundefinedtheundefinedtwoundefinedlimitsundefinedbelowundefinedshouldundefinedbeundefinedequalundefinedatundefined$0\in \mathbb{C}$:
\begin{equation}
\label{complex_derivative_limit}
\lim_{\Delta x\to 0}\frac{f(x_0 + \Delta x + y_0i) -f(x_0 + y_0 i)}{\Delta x} = \lim_{\Delta y\to 0}\frac{f(x_0 + (y_0 +\Delta y)) -f(x_0 + y_0 i)}{i\Delta y }
\end{equation}
Butundefinedcomputingundefinedtheundefinedlimitundefinedtellsundefinedusundefinedthat:
\begin{equation}
2x_0 \neq 1
\end{equation}
Unlessundefined$x_0 = \frac{1}{2}$.undefinedHenceundefinedtheundefinedfunctionundefinedisundefinedonlyundefineddifferentiableundefinedonundefinedtheundefinedlineundefined$x_0=\frac12$.

Oneundefinedmightundefinedask,undefineddoesn'tundefinedtheundefinedseemundefinedproblemundefinedexistsundefinedforundefinedtheundefinedrealundefinedanalog?undefinedIt'sundefinedquiteundefinedobviousundefinedtheundefinedanswerundefinedisundefinedno,undefinedsinceundefinedweundefinedneverundefinedcomputeundefinedaundefinedsingleundefinedlimitundefinedofundefinedaundefinedfiniteundefineddifference.

Inundefinedaundefinedway,undefined$\mathbb{C} \to \mathbb{C}$undefinedisundefinedlikeundefined$\mathbb{R}\to \mathbb{R}$,undefinedweundefineddefineundefinedtheundefinedderivativeundefinedtoundefinedbeundefinedtheundefinedlimitundefinedofundefinedfiniteundefineddifference.undefinedYetundefinedtheundefinedlimits

\begin{equation}
\lim_{\Delta z\to 0}\frac{u(z+\Delta z) -u(z)}{\Delta z}, \lim_{\Delta z\to 0}\frac{v(z+\Delta z) -v(z)}{\Delta z}
\end{equation}

resemblesundefinedtheundefined$\mathbb{R}^2 \to \mathbb{R}$undefinedwhereundefinedweundefinedexpectundefinedtheundefinedresultundefinedtoundefinedbeundefinedindependentundefinedofundefineddirection.undefinedThat'sundefinedwhyundefinedtheundefinedrealundefinedanalogundefinedisundefinedfundamentallyundefineddifferentundefinedfromundefinedtheundefinedcomplexundefinedfunctionundefineditself.

\end{mystthmexample}

\subsubsection{Cauchy-RiemannundefinedCondition}

Expandingundefined(\ref{complex_derivative_limit}),undefinedweundefinedwouldundefinedget:
\begin{equation}
\begin{aligned}
\lim_{\Delta x\to 0} \frac{f(x_0+\Delta x + y_0i) - f(x_0+y_0i)}{\Delta x} &= \lim_{\Delta y\to 0} \frac{f(x_0 + (y_0+\Delta yi)) - f(x_0+y_0i)}{\Delta y} \\
\lim_{\Delta x\to 0}\frac{u(x_0+\Delta x, y_0)-u(x_0,y_0)}{\Delta x} + i \frac{v(x_0+\Delta x, y_0)-v(x_0,y_0)}{\Delta x} &=  
\lim_{\Delta y\to 0}\frac{u(x_0, y_0+\Delta y)-u(x_0,y_0)}{i\Delta y} + i \frac{v(x_0, y_0+\Delta y)-v(x_0,y_0)}{i\Delta y}\\
\frac{\partial u}{\partial x}+i\frac{\partial v}{\partial x} &= -\frac{\partial u}{\partial y}i + \frac{\partial v}{\partial y}
\end{aligned}
\end{equation}
Noteundefinedthatundefinedthisundefinedequationundefinedessentiallyundefinedunwindundefinedtheundefined``mixingundefinedofundefinedcomponents''undefinedweundefinedmentionedundefinedinundefinedDifferenceundefined2undefined:undefinedMixingundefinedOfundefinedComponents.undefinedHenceundefinedweundefinedget:

\begin{equation}
\label{cr}
\boxed{\frac{\partial u}{\partial x} = \frac{\partial v}{\partial y} \qquad 
\frac{\partial v}{\partial x} =  - \frac{\partial u}{\partial y}}
\end{equation}

ThisundefinedisundefinedknownundefinedasundefinedtheundefinedCauchy-RiemannundefinedCondition(C-Rundefinedcondition).

Let'sundefinedrecallundefinedwhatundefinedhadundefinedhappenedundefinedsoundefinedfar,undefinedweundefineddiscussedundefinedtheundefinedideaundefinedofundefinedcomplexundefinedderivativeundefinedbeingundefinedaundefinedlimitundefinedofundefinedfiniteundefineddifferenceundefinedleadingundefinedtoundefinedcomputingundefinedlimitsundefinedofundefinedtheundefinedformundefined$\mathbb{R}^2\to \mathbb{R}$(twoundefinedinput,undefinedoneundefinedoutput)undefinedwhichundefinedputsundefinedadditionalundefinedrestrictionundefinedonundefineddifferentiability.undefinedToundefinedbeundefinedmoreundefinedspecific,undefinedinundefinedorderundefinedforundefinedtheseundefinedlimitundefinedtoundefinedexists,undefinedtheundefinedvalueundefinedofundefinedtheundefinedlimitundefinedmustundefinedbeundefinedinvariantundefinedofundefineddirectionsundefinedapproachingundefinedtheundefinedpointundefinedofundefinedinterestundefined$x_0, y_0$.undefinedConsideringundefinedtheundefinedspecificundefineddirectionundefinedalongundefinedtheundefinedrealundefinedandundefinedcomplexundefinedaxisundefinedweundefinedgetundefinedtheundefinedso-calledundefinedCRundefinedcondition.

WeundefinedhaveundefinedeffectivelyundefinedshownundefinedthatundefinedCRundefinedconditionundefinedisundefined\textit{necessary}undefinedforundefinedcomplexundefineddifferentiability.undefinedItundefinedturnsundefinedoutundefinedhowever,undefinedthatundefinedCRundefinedconditionundefinedisundefinedalsoundefined\textit{sufficient}.

\begin{mystthmtheorem}{Cauchy-Riemann}{theorem-analyticalfunction-4}Aundefinedcomplexundefinedfunctionundefined$f:\mathbb{C} \to \mathbb{C}$undefinedisundefineddifferentiableundefinedatundefined$z_0=x_0+iy_0$undefinediffundefineditundefinedsatisfiesundefinedtheundefined\textit{Cauchy-Riemannundefinedcondition}:

\begin{enumerate}
\item $u,v\in C^1$(continouslyundefineddifferentiable)
\item Theundefinedpartialundefinedderivativesundefinedsatisfy
\begin{equation}
\begin{aligned}
\frac{\partial u}{\partial x} &= \frac{\partial v}{\partial y} \\
\frac{\partial u}{\partial y} &= -\frac{\partial v}{\partial x}
\end{aligned}
\end{equation}
Inundefinedwhichundefinedcase,undefinedtheundefinedderivativeundefinedis:
\begin{equation}
\begin{aligned}
\frac{\partial f}{\partial z} &= \frac{\partial u}{\partial x} + i \frac{\partial v}{\partial x} \\
&= \frac{\partial v}{\partial y} - i \frac{\partial u}{\partial y}
\end{aligned}
\end{equation}
\end{enumerate}

\end{mystthmtheorem}
\begin{mystthmproof}{CRundefinedconditionundefinedisundefinedsufficient}{proof-analyticalfunction-5}Weundefinedneedundefinedtoundefinedshow,undefined$\forall \epsilon >0$,undefinedundefined$\exists \delta >0$undefinedsuchundefinedthatundefined$|\Delta z|<\delta$undefinedimpliesundefined:
\begin{equation}
\left| \frac{f(z_0+\Delta z)-f(z_0)}{\Delta z} - \left(\frac{\partial u}{\partial x} + i\frac{\partial v}{\partial x}\right) \right| < \epsilon
\end{equation}
Weundefinedconsiderundefinedtheundefinedlinaerundefinedapproximationundefinedofundefined$f(z_0+\Delta z) -f(z_0)$undefinedaroundundefined$z_0$:
\begin{equation}
\begin{aligned}
f(z_0+\Delta z)-f(z_0)&=\textcolor{blue}{u(x_0+\Delta x,y_0+\Delta y)-u(x_0,y_0)} + i[\textcolor{red}{v(x_0+\Delta x, y_0+\Delta y)-v(x_0,y_0)}]
\\ &+ O(\Delta x^2) + iO(\Delta y^2) \\
&= \textcolor{blue}{\frac{\partial u}{\partial x}\Delta x +\frac{\partial u}{\partial y} \Delta y} +i\left(\textcolor{red}{\frac{\partial v}{\partial x}\Delta x +\frac{\partial v}{\partial y}\Delta y} \right) \\ &+ O(\Delta x^2)+iO(\Delta y^2) \\
&= \left(\textcolor{blue}{\frac{\partial u}{\partial x}}+i\textcolor{red}{\frac{\partial v}{\partial x}}\right)\Delta x+\left(\textcolor{blue}{\frac{\partial u}{\partial y}} + i\textcolor{red}{\frac{\partial v}{\partial y}}\right)\Delta y\\ &+ O(\Delta x^2)+iO(\Delta y^2) \\
\end{aligned}
\end{equation}
Weundefinedneedundefinedtoundefinedexposeundefined$i\Delta y$,undefinedhenceundefinedweundefinedrewriteundefinedtheundefinedthirdundefinedtermundefinedinundefinedtheundefinedlastundefinedequationundefinedabove:
\begin{equation}
\left(\textcolor{blue}{\frac{\partial u}{\partial y}} + i\textcolor{red}{\frac{\partial v}{\partial y}}\right)\Delta y \longrightarrow \left(\textcolor{blue}{-\frac{\partial u}{\partial y}}i + \textcolor{red}{\frac{\partial v}{\partial y}}\right)i\Delta y
\end{equation}
Therefore:
\begin{equation}
\begin{aligned}
f(z_0+\Delta z)-f(z_0) &=  \left(\textcolor{blue}{\frac{\partial u}{\partial x}}+i\textcolor{red}{\frac{\partial v}{\partial x}}\right)\Delta x+\left(\textcolor{blue}{-\frac{\partial u}{\partial y}}i + \textcolor{red}{\frac{\partial v}{\partial y}}\right)i\Delta y\\ &+ O(\Delta x^2)+iO(\Delta y^2)
\end{aligned}
\end{equation}
NowundefinedapplyingundefinedCRundefinedcondition:
\begin{equation}
\begin{aligned}
f(z_0+\Delta z)-f(z_0) &=  \left(\textcolor{blue}{\frac{\partial u}{\partial x}}+i\textcolor{red}{\frac{\partial v}{\partial x}}\right)\Delta x+\left(\textcolor{red}{\frac{\partial v}{\partial x}}i + \textcolor{blue}{\frac{\partial u}{\partial x}}\right)i\Delta y\\ &+ O(\Delta x^2)+iO(\Delta y^2)  \\
&= \textcolor{blue}{\frac{\partial u}{\partial x}}\left(\Delta x + i\Delta y\right) + i\textcolor{red}{\frac{\partial v}{\partial x}}(\Delta x+i\Delta y)
\\ &+ O(\Delta x^2)+iO(\Delta y^2) \\
&= \left(\textcolor{blue}{\frac{\partial u}{\partial x}} + i\textcolor{red}{\frac{\partial v}{\partial x}}\right)\underbrace{(\Delta x +i\Delta y)}_{\Delta z}\\
&+ \underbrace{O(\Delta x^2)+iO(\Delta y^2)}_{O(|\Delta z|\Delta z)} \\
\end{aligned}
\end{equation}
Thereforeundefinedtheundefinedquotientundefinedis:
\begin{equation}
\frac{f(z_0+\Delta z)-f(z_0)}{\Delta z} =\textcolor{blue}{\frac{\partial u}{\partial x}} + i\textcolor{red}{\frac{\partial v}{\partial x} }+O(|\Delta z|)
\end{equation}
andundefinedthusundefinedtheundefinedquanityundefinedweundefinedwantundefinedtoundefinedcontrolundefinedis:
\begin{equation}
\left| \frac{f(z_0+\Delta z)-f(z_0)}{\Delta z} - \left(\frac{\partial u}{\partial x} + i\frac{\partial v}{\partial x}\right) \right|=O(|\Delta z|)
\end{equation}
Byundefineddefinitionundefinedofundefinedbig-O,undefined$O(|\Delta z|)<C|\Delta z|$undefinedforundefined$|\Delta z|\in \mathbb{R}$.undefinedHenceundefined$\forall \epsilon >0$,undefinedweundefinedcanundefinedpickundefined$\delta =\frac{\epsilon}{C}$undefinedsuchundefinedthatundefined$O(|\Delta z|)<C\delta =\epsilon$

\end{mystthmproof}
\subsubsection{AnalyticalundefinedFunction}

Nowundefinedweundefinedareundefinedreadyundefinedforundefinedtheundefineddefinitionundefinedofundefined\textbf{analyticundefinedfunction}.

\begin{mystthmdefinition}{AnalyticundefinedFunctionundefinedonundefinedEntireundefinedFunction}{definition-analyticalfunction-6}\begin{itemize}
\item Aundefinedcomplexundefinedfunctionundefined$f:\mathbb{C} \to \mathbb{C}$undefinedisundefined\textbf{analyticundefinedatundefined$z_0$}undefinedifundefineditundefinedisundefineddifferentiableundefinedatundefined$z$undefinedandundefinedaundefinedneighborhoodundefinedofundefined$z$.
\item Aundefinedfunctionundefinedthatundefinedisundefineddifferentiableundefinedthourhgoutundefinedsomeundefinedopenundefinedsetundefined$\Omega \in \mathbb{C}$undefinedisundefined\textbf{analyticundefinedinundefined$\Omega$}
\item Aundefinedfunctionundefinedthatundefinedisundefinedanalyticalundefinedthroughoutundefined$\mathbb{C}$undefinedisundefinedcalledundefined\textbf{entire}.
\end{itemize}

\end{mystthmdefinition}
Analyticundefinedisundefinedaundefinedmuchundefinedmoreundefinedrestrictiveundefinedpropertyundefinedthanundefinedmerelyundefineddifferentiable(sinceundefineditundefinedrequiresundefineddifferentiabilityundefinedinundefinedtheundefinedneighborhood).undefinedForundefinedinstance,undefinedweundefinedcanundefinedeasilyundefinedconstructundefinedcomplexundefinedfunctionsundefinedthatundefinedisundefineddifferentiableundefinedonundefined\textit{infinitelyundefinedmanyundefinedpoints}undefinedinundefined$\mathbb{C}$undefinedbutundefinedanalyticundefined\textit{nowhere}.

\begin{mystthmexample}{Differentiableundefinedonundefinedinfiniteundefinedmanyundefinedpointsundefinedbutundefinedanalyticundefinednowhere}{example-analyticalfunction-7}Supposeundefinedweundefinedwantundefinedtheundefinedfunctionundefinedtoundefinedbeundefineddifferentiableundefinedonundefined$y=x$undefinedandundefined$y= -x$.undefinedThenundefinedweundefinedcan,undefinedbyundefinedCR,undefinedlet:
\begin{equation}
\frac{\partial u}{\partial x} =2x, \frac{\partial v}{\partial y} = 2y
\end{equation}
and
\begin{equation}
\frac{\partial v}{\partial x} =-2x, \frac{\partial u}{\partial y} = 2y
\end{equation}
Aundefinedeasyundefinedchoiceundefinedwouldundefinedjustundefinedbe:
\begin{equation}
\begin{aligned}
f&=x^2+y^2 +i(y^2-x^2) \\
&= zz^*+i\left(\frac{z^2+{z^*}^2}{2}\right)
\end{aligned}
\end{equation}

\end{mystthmexample}
\begin{mystadmonition}{mystAmber}{QuickundefinedExercise}\label{ananowhere}Doundefinedyouundefinedseeundefinedwhyundefinedtheundefinedaboveundefinedfunctionundefinedisundefinedanalyticalundefinednowhere?

\end{mystadmonition}
\begin{mystadmonition}{mystGreen}{Solutionundefinedtoundefined\ref{ananowhere}}Everyundefinedopenundefinedsetundefinedinundefined$\mathbb{C}$undefinedcontainsundefined\textit{some}undefinedpointundefinednotundefinedonundefined$y=x$undefinedorundefined$y= -x$

\end{mystadmonition}